% The formats reference: the data model of a recording (viewers/vitalstream/docs/data-model.md),
% the index format, version 2 (viewers/vitalstream/docs/index-format.md), and the files of a WFDB
% collection (docs/wfdb-collection.md: index.csv, gaps.csv, vitals.csv, events.csv, markers),
% and the reading rules of EDF, EDF+, BDF and BDF+ (docs/edf.md, "Reading rules"; checked against
% libs/engine/src/sources/edf.rs). How records, collections and folders are read is
% chap:recordings.
\chapter{Formats reference}
\label{chap:formats-reference}

This chapter specifies the data that Medorama and the VitalStream viewers share. It is written
for implementers of readers, adapters and converters, and for anyone who checks them. It covers
four subjects:
\begin{itemize}
\item the \textbf{data model} of a recording, which every reader builds whatever the recording's
  format;
\item the \textbf{index}, the files that let the viewers show days of waveforms at once, which
  Medorama, the three viewers and the adapters write byte for byte alike;
\item the files of a \textbf{WFDB collection}, the folder of WFDB records with a few CSV files that
  converters write;
\item the \textbf{reading rules of EDF}, EDF+, BDF and BDF+, which the three engines implement
  identically.
\end{itemize}
How a single record, a multi-segment record, a collection, a folder of records or an EDF file is
read and placed on the time axis is described for users in \cref{chap:recordings}; how to add a
reader for a new format is described in \cref{chap:readers}.

\section{The data model}
\label{sec:formats-reference-model}

The data model describes, in language-neutral terms, what the viewers make of a recording,
whatever its format. A reader (the WFDB or EDF reader of the engines, or an adapter,
\cref{chap:addons}) builds it from the recording, the index stores it
(\cref{sec:formats-reference-index}), and the viewer API serves it (\cref{chap:viewer-api}).

\subsection{Principles}
\label{sec:formats-reference-principles}

\begin{itemize}
\item \textbf{Only what the recording states.} The readers do no resampling, interpolation, gap
  filling, re-timing or deduplication: every sample drawn is a stored sample.
\item \textbf{Ambiguity is shown, not resolved.} Where the recording does not settle something (a
  break's exact length, a start time), the index says so, with \code{exact}, \code{timed} and
  notes.
\item \textbf{Times are UTC}, as Unix seconds. The display time zone is a setting of the viewer
  only.
\end{itemize}

\subsection{Structure}
\label{sec:formats-reference-structure}

\begin{xltabular}{\linewidth}{@{}P{2.2cm}L@{}}
\caption{The terms of a recording's structure.}
\label{tab:formats-reference-structure}\\
\toprule
\tabhead Term & \tabhead Meaning \\
\midrule
\endfirsthead
\toprule
\tabhead Term & \tabhead Meaning \\
\midrule
\endhead
\bottomrule
\endfoot
channel & one signal as the user sees it (ECG lead II, ART, Pleth, \dots): a label, a unit, a
  modality \\
piece & samples of one channel that the recording proves contiguous, at one constant sampling
  rate \code{fs}: sample $i$ is at $t_0 + i / f_s$. A channel is the pieces it is made of; pieces
  are drawn apart, never joined \\
lane & the pieces of a channel never overlap in time. In WFDB records, a piece that would overlap
  another of its channel (a signal recorded twice at the same time) goes into the first lane of
  that channel it does not overlap; each further lane is a channel of its own, labelled with its
  number: \code{II}, \code{II (2)}, \code{II (3)} \\
timed & false when the recording states no time for a piece: it is placed next to its neighbors
  and marked \ui{time not stated} \\
envelope & per piece, the minimum and maximum of every block of 32 samples, and of every 4 blocks
  above that: what long windows are drawn from \\
break & a place where contiguity is not proven: its reason, its length when known
  (\code{gap\_s}), whether that length is \code{exact}, how it was measured, and the channels it
  concerns \\
\end{xltabular}

\subsection{Values}
\label{sec:formats-reference-values}

\begin{itemize}
\item \textbf{Units:} physical units where the recording states a calibration; \code{adu} where it
  states none (the value drawn is the stored count).
\item \textbf{Nothing is assumed} that a recording does not state: a WFDB record whose header
  gives no sampling frequency is not read (opened alone, it is refused; in a folder of plain
  records or a WFDB collection it is left out, listed in the index's \code{skipped} with the
  reason and counted in the notes; a collection of which no record can be read is refused), a
  signal with a gain of 0 is shown in
  stored counts (\code{adu}), and a signal whose header gives no units has none. WFDB's old
  defaults of 250\,Hz, a gain of 200 and millivolts are not used.
\item \textbf{Missing samples} are NaN, drawn as a gap in the line, never as zero. In WFDB they
  are the format's invalid value.
\item \textbf{Zero is a value}, never ``missing''.
\end{itemize}

\subsection{Vital signs, alarms and markers}
\label{sec:formats-reference-events}

\begin{itemize}
\item \textbf{A vital-sign series:} values at their times, by name and unit; one series per name
  and source, never merged (\cref{sec:formats-reference-vitals}).
\item \textbf{An alarm:} an interval (start, end) with a text, a kind and a level: 1 advisory,
  2 warning, 3 crisis. Alarm colors are the only red on the viewers' pages.
\item \textbf{A marker:} a mark at a sample of a piece (an R wave, a pacing spike), by name.
\end{itemize}

\subsection{Identifiers}
\label{sec:formats-reference-identifiers}

The viewers never read, keep or show a patient's identity: the index holds no name, identifier or
date of birth, and neither do the viewers' answers. A recording's path and file name may still
identify a patient (bedside-monitor systems name files by unit, bed and time), and every value of
a recording is patient data: keep the recordings, and the index folders made from them, on
storage protected like the recordings.

\section{The index}
\label{sec:formats-reference-index}

The index of a recording holds its structure (its channels, and the pieces each channel is made
of), the min/max envelopes of every waveform, and its markers, vital-sign series, alarms and
breaks. This section specifies the index, version 2, so that programs written independently write
and read the same files and can share one index folder: the viewers \cmd{vitalstream-rs} (with
the Rust engine), \cmd{vitalstream-cpp} (with the C++ engine) and \cmd{vitalstream-py}, Medorama
(with the Rust engine), and the adapters. All of them write the same bytes for the same
recording.

\subsection{The index folder}
\label{sec:formats-reference-folder}

An index folder holds one subfolder per indexed recording, named by its key
(\cref{sec:formats-reference-key}):
\begin{lstlisting}[style=text]
<index folder>/
    0f38a5788406244247a4/            one recording
        meta.json
        env_min_0.f32  env_max_0.f32  env_min_1.f32 ...
        mk_piece.i32  mk_sample.i64  mk_code.i16
        vit_t.f64  vit_v.f64
    ...
\end{lstlisting}
Any other entry of an index folder (temporary folders, the Python viewer's progress files
\file{<key>.progress.json}) is ignored by readers. Medorama's index folder is its
\code{index\_dir}, a viewer's its \code{cache\_dir} (\cref{chap:settings-reference}).

\subsection{The key}
\label{sec:formats-reference-key}

The key identifies a recording \textbf{as it is now}: it changes when the file is changed, moved
or replaced, so that an outdated index is never used.
\begin{lstlisting}[style=text]
key = the first 20 hex digits (lowercase) of SHA-256( JSON text of [version, path, fingerprint] )
\end{lstlisting}

\begin{xltabular}{\linewidth}{@{}P{3.5cm}L@{}}
\caption{The elements of the key.}
\label{tab:formats-reference-key}\\
\toprule
\tabhead Element & \tabhead Value \\
\midrule
\endfirsthead
\toprule
\tabhead Element & \tabhead Value \\
\midrule
\endhead
\bottomrule
\endfoot
\code{version} & the index format's version, \code{2} \\
\code{path} & the recording's absolute path with every symbolic link resolved (the real path). On
  Windows, without the \code{\textbackslash\textbackslash?\textbackslash} prefix:
  \code{C:\textbackslash data\textbackslash x.Stp}, and
  \code{\textbackslash\textbackslash server\textbackslash share\textbackslash x.Stp} for network
  shares \\
\code{fingerprint}, a file & \code{[size, mtime]}: the size in bytes; the modification time in
  \textbf{integer nanoseconds} since 1970-01-01 UTC (on Windows, from the file's last-write time,
  which has a resolution of 100\,ns) \\
\code{fingerprint}, a folder (WFDB) & the list of \code{[name, size, mtime]} of \textbf{every
  file} below the folder (subfolders included; folders themselves and hidden entries, whose names
  start with \code{.}, excluded; a symbolic link to a file counted under the link's name, with its
  target's size and mtime; links to folders and broken links excluded), \code{name} relative to
  the folder with \code{/} between its components on every system, sorted by \code{name} (then
  size, mtime). The folder is walked depth first, the entries of each folder in the byte order of
  their UTF-8 names, down to \code{max\_depth} levels (20 by default), and the walk stops once
  \textbf{200\,000} files are listed \\
\code{fingerprint}, a WFDB record (a \file{.hea} file) & the same list for \textbf{the record's
  files}: its header, the signal files its header names (other than \code{-} and
  \code{\textasciitilde}), its \file{.mrk} file if there is one, and, for a multi-segment record,
  the same for each segment and the layout header; \code{name} relative to the header's folder,
  each file once, sorted by \code{name} \\
\end{xltabular}

\paragraph{The JSON text} is compact, with no spaces: \code{[2,"/data/x.Stp",[133325312,1700000000123456789]]}.
Strings are in double quotes, with the escapes \code{\textbackslash"}, \code{\textbackslash\textbackslash},
\code{\textbackslash n}, \code{\textbackslash r}, \code{\textbackslash t}, \code{\textbackslash b},
\code{\textbackslash f}, and every other character below U+0020 or above U+007E as
\code{\textbackslash u} and 4 lowercase hex digits (UTF-16 surrogate pairs above U+FFFF);
integers are in decimal.

\paragraph{An example} (a made-up file): for \file{/data/rec.Stp}, 1000 bytes, modified at
1700000000.5\,s, the text is
\begin{lstlisting}[style=text]
[2,"/data/rec.Stp",[1000,1700000000500000000]]
\end{lstlisting}
and the key, the first 20 hex digits of that text's SHA-256, is \code{f61f534c1e1a4d244d7a}.

\subsection{The files of one index}
\label{sec:formats-reference-files}

\begin{xltabular}{\linewidth}{@{}P{5.0cm}L@{}}
\caption{The files of one index.}
\label{tab:formats-reference-files}\\
\toprule
\tabhead File & \tabhead Content \\
\midrule
\endfirsthead
\toprule
\tabhead File & \tabhead Content \\
\midrule
\endhead
\bottomrule
\endfoot
\file{meta.json} & the recording's structure (\cref{sec:formats-reference-meta}), UTF-8 JSON \\
\file{env\_min\_<L>.f32}, \file{env\_max\_<L>.f32} & envelope level \code{L} = 0, 1, \dots:
  the minima and maxima of all pieces, concatenated \\
\file{mk\_piece.i32}, \file{mk\_sample.i64}, \file{mk\_code.i16} & markers: piece, sample number,
  index into \code{marker\_names} \\
\file{vit\_t.f64}, \file{vit\_v.f64} & vital-sign times (Unix seconds, UTC) and values, all series
  concatenated \\
\end{xltabular}

\paragraph{Array files} hold raw little-endian values of the type their extension names
(\code{f32}, \code{f64}, \code{i16}, \code{i32}, \code{i64}), with no header; the number of values
is the file's size divided by the value's size, and an array with no values is an empty file. The
marker and vital-sign files always exist; envelope files exist for the levels 0 to $L-1$ only.

\subsection{meta.json}
\label{sec:formats-reference-meta}

A JSON object, written compact (no spaces) and in ASCII, with the escapes of the key's text.
Integers are written without a decimal point; other numbers in the shortest form that reads back
exactly (\code{0.1}, \code{1e-05}, \code{1709287200.0}); NaN and infinities never appear: they are
written as \code{null}.

\begin{xltabular}{\linewidth}{@{}P{3.0cm}P{2.2cm}L@{}}
\caption{The keys of \file{meta.json}.}
\label{tab:formats-reference-meta}\\
\toprule
\tabhead Key & \tabhead Type & \tabhead Meaning \\
\midrule
\endfirsthead
\toprule
\tabhead Key & \tabhead Type & \tabhead Meaning \\
\midrule
\endhead
\bottomrule
\endfoot
\code{index\_version} & integer & \code{2} \\
\code{key} & string & the key (\cref{sec:formats-reference-key}) \\
\code{path} & string & the real path used in the key \\
\code{name} & string & the file or folder name as opened \\
\code{kind} & string & the recording kind: \code{"wfdb"}, \code{"edf"}, or the kind of the adapter
  that read it (\code{"stp"}) \\
\code{format} & string & the format's label: \code{"WFDB record"}, \code{"WFDB multi-segment
  record"}, \code{"WFDB collection"}, \code{"WFDB records"} (a folder of plain records),
  \code{"EDF"}, \code{"EDF+"}, \code{"BDF"}, \code{"BDF+"}, or the adapter's
  (\code{"BedMasterEx .Stp"}) \\
\code{adapter} & string & the adapter that read the recording, when one did; it reads its samples
  again on demand \\
\code{description} & string & a sentence about the recording from its adapter, optional \\
\code{engine} & string & the version of the program that wrote the index (informative) \\
\code{indexed\_s} & number & the seconds the indexing took (informative) \\
\code{time\_resolution\_s} & number or null & the time resolution of the source's time stamps \\
\code{channels} & list & per channel: \code{id} (its position), \code{label}, \code{unit},
  \code{modality}, and source-specific fields \\
\code{pieces} & list & per piece (a stretch of one channel the source proves contiguous):
  \code{id} (its position), \code{ch} (the channel's id), \code{fs} (Hz), \code{n} (samples),
  \code{t0} (Unix seconds, UTC, of sample 0), \code{timed} (false when the source states no time;
  omitted: true), \code{env} (per level, \code{[offset, length]} into that level's envelope
  arrays), and source-specific fields (a WFDB piece's \code{record} and \code{sig}; an EDF
  piece's \code{sig} and \code{data\_record}) \\
\code{vitals} & list & per series: \code{name}, \code{unit}, \code{offset}, \code{n} into
  \file{vit\_t} and \file{vit\_v} \\
\code{alarms} & list & \code{t0}, \code{t1}, \code{text}, \code{level}, \code{kind} \\
\code{gaps} & list & the breaks of the recording: \code{t0}, \code{t1}, \code{reason},
  \code{gap\_s}, \code{exact}, \code{measured\_by}, \code{notes}, optional \code{channels} \\
\code{untimed} & list & the intervals without a stated time: \code{t0}, \code{t1} \\
\code{late\_blocks} & list & optional: data blocks the source stored out of order: \code{t0},
  \code{t1} (the block's time span; \code{t1} null when unknown), \code{moved} (true: put back in
  order, so no break; false: left in file order, and the break it causes is in \code{gaps} with a
  note), \code{reason} (why not moved; empty when moved), \code{late\_by\_s},
  \code{late\_by\_blocks} (null when not counted), \code{segment}. The viewers mark the moved ones
  only with their \code{qc\_verbose} setting \\
\code{marker\_names} & list of strings & the marker names, indexed by \file{mk\_code} \\
\code{skipped} & object & what could not be shown: an adapter's counts; for WFDB records, the
  records left out, by name, with the reason \\
\code{notes} & list of strings & remarks for the user \\
\code{reader} & object & what the source's reader needs to read samples again (an adapter's
  state); written and read only by the program that reads the source, and no other may depend on
  it \\
other keys & & source-specific; readers ignore the keys they do not know \\
\end{xltabular}

\paragraph{The order of the keys.} The reader's keys come first, in the order it writes them, then
\code{kind}, \code{format} and \code{adapter}, then \code{path}, \code{name}, \code{key},
\code{index\_version}, \code{engine}, \code{indexed\_s}. Programs that share an index folder write
the same order, so that they write the same bytes.

\subsection{Arrays}
\label{sec:formats-reference-arrays}

\paragraph{Envelopes.} Level 0 holds, for each piece, the minimum and the maximum of every block of
\textbf{32} samples, in order; the last block may be shorter. Each further level combines
\textbf{4} consecutive buckets of the level below, until a piece has a single bucket. Minima and
maxima ignore NaN (NumPy's \code{fmin} and \code{fmax}); a bucket whose samples are all NaN is
NaN. Values are in the channel's physical units. Piece $p$'s buckets at level $L$ are
\code{env\_min\_L[o : o+len]} and \code{env\_max\_L[o : o+len]}, where \code{[o, len]} is
\code{pieces[p].env[L]}.

\paragraph{Markers} are sorted by (piece, sample); markers with equal values keep the order in
which the source gave them. A marker's time is \code{t0 + sample / fs} of its piece.

\paragraph{Vital signs.} Series are ordered by their number of values (most first), then by name;
within a series, values are sorted by time; only finite times and values are stored. Series $s$ is
\code{vit\_t[offset : offset+n]} and \code{vit\_v[offset : offset+n]}.

\subsection{Writing and sharing}
\label{sec:formats-reference-writing}

\begin{itemize}
\item \textbf{Atomic writing:} a writer builds the index in a temporary folder next to the final
  one, named uniquely for that writer (\file{<key>.<process id>.<n>.tmp}, \code{n} counting the
  writer's builds), then renames it to \file{<key>}. The temporary folder is removed on every way
  out (an error, a cancellation, a failure of the program).
\item \textbf{One writer per key:} within a program, builds of the same key wait for each other.
\item \textbf{Keeping or replacing:} when the build is complete, the writer checks the folder
  \file{<key>}. If it holds a complete index (\cref{sec:formats-reference-loading}) written by
  the same program version (\code{engine}), it is kept and the new build is discarded. Otherwise
  the old folder, if any, is first renamed aside (\file{<key>.<process id>.<n>.old}), the new one
  renamed to \file{<key>}, and the old one removed. Readers therefore see a complete index, or
  none for the moment between the two renames.
\item \textbf{Shared folders:} several programs (Medorama and the three viewers) may use one index
  folder at the same time; when two write the same key, the content is identical, and either one
  is kept.
\item \textbf{Read only after writing:} an index is never modified after it is written, so that
  readers may memory-map its arrays.
\item \textbf{Version check:} a reader uses an index only if \file{meta.json} exists, its
  \code{index\_version} equals the version it implements, and the index passes the checks of
  \cref{sec:formats-reference-loading}; otherwise it indexes again.
\item \textbf{Never next to the recording:} an index folder is never inside a folder of
  recordings.
\end{itemize}

\subsection{Loading an index}
\label{sec:formats-reference-loading}

A reader checks an index before it uses it, so that no query reads outside an array, whoever
wrote the files. An index that fails any check is treated as not indexed. A writer runs the same
checks on what it has just built, and refuses to keep an index that fails them (for example times
or rates that are not finite numbers, which \file{meta.json} would hold as \code{null}).
\begin{enumerate}
\item \file{meta.json} is a JSON object whose \code{index\_version} is the reader's version.
\item Every array file exists (\code{env\_max\_L} for every \code{env\_min\_L} found, the levels
  counted from 0 until \code{env\_min\_L} is missing), and its size is a whole number of values.
\item \code{channels}, \code{pieces}, \code{vitals} and \code{marker\_names} are lists when
  present (missing: empty).
\item Piece \code{i}: \code{id} is the integer \code{i}; \code{ch} is an integer with
  $0 \le$ \code{ch} $<$ the number of channels; \code{fs} is a finite number greater than 0;
  \code{t0} a finite number; \code{n} an integer, 0 or more; \code{env}, when present, a list of
  at most as many entries as there are envelope levels, entry \code{L} being \code{[o, len]}, two
  integers of 0 or more with \code{o + len} at most the number of values of the shorter of
  \code{env\_min\_L} and \code{env\_max\_L}.
\item \code{mk\_piece}, \code{mk\_sample} and \code{mk\_code} have the same number of values;
  every \code{mk\_piece} value is a piece id and the values never decrease; every \code{mk\_code}
  value is an index into \code{marker\_names}.
\item \code{vit\_t} and \code{vit\_v} have the same number of values, and every series has
  integers \code{offset} and \code{n}, 0 or more, with \code{offset + n} at most that number.
\end{enumerate}

\subsection{Limits}
\label{sec:formats-reference-limits}

\begin{itemize}
\item \textbf{Pieces.} A recording has at most \textbf{100\,000} pieces. A reader that would create
  more (an EDF+D file with many short runs, a folder of many records with many signals) refuses
  the recording with \code{the recording has more than 100000 pieces (stretches of one signal):
  it is not indexed}. The check is made before each new run (EDF) or record part (WFDB) adds its
  pieces.
\item \textbf{Marker names.} \code{mk\_code} is 16-bit, so at most \textbf{32\,768} distinct
  names are kept. The first 32\,767 names keep codes 0 to 32\,766; the next name found is written
  as \code{(other markers)} with code 32\,767, and every later new name gets that code too. The
  WFDB reader then adds the note \code{More than 32767 distinct marker names: the others are shown
  as "(other markers)".}
\end{itemize}

\subsection{Files outside the data roots}
\label{sec:formats-reference-roots}

A program that limits the folders it opens (the viewers' data roots) also limits every file a
reader opens for a recording, symbolic links resolved: a WFDB record's header, the signal files
its header names (an absolute name as given, a name starting with \code{\textasciitilde} in the
home folder, any other next to the header), its \file{.mrk} file, each segment's header and
signal files, and the CSV files of a folder or collection. A file that does not exist is allowed
(nothing is read from it); a broken link is not. This is checked before a record is indexed and
again each time its samples are read:
\begin{itemize}
\item one record (a \file{.hea} file opened): refused (\code{403} from the viewer's \code{open});
\item a record of a folder or collection: left out, in \code{skipped}, with the reason
  \code{<file> lies outside the data roots: the record <name> is not read};
\item samples read later from a file that now leads outside: an error.
\end{itemize}
The Rust engine does this once a program calls \code{recording\_engine::paths::set\_allowed\_roots}.

\subsection{Versions and conformance}
\label{sec:formats-reference-versions}

\begin{xltabular}{\linewidth}{@{}P{1.6cm}L@{}}
\caption{The versions of the index format.}
\label{tab:formats-reference-versions}\\
\toprule
\tabhead Version & \tabhead Changes \\
\midrule
\endfirsthead
\toprule
\tabhead Version & \tabhead Changes \\
\midrule
\endhead
\bottomrule
\endfoot
1 & the first version \\
2 & this specification: one WFDB record (a \file{.hea} file) and its key; multi-segment records;
  folders of plain records read with their subfolders (hidden entries and symbolic links skipped,
  in the folder's key too); overlapping pieces of a channel in lanes of their own. Later within
  version 2, with no existing index changed: EDF, EDF+, BDF and BDF+ files (\code{kind}
  \code{edf}); a folder's key counts the files its symbolic links lead to, and its walk is
  bounded (the keys of folders without links, of at most 200\,000 files within \code{max\_depth}
  levels, are unchanged); a WFDB header that states more samples than its signal files hold is
  read to the end of its data; the limits and the checks on loading \\
\end{xltabular}

Any change to the files, to \file{meta.json} or to the key makes a new version, and indexes of
other versions are built again. The script \file{viewers/vitalstream/tests/parity/index\_parity.py}
indexes recordings with every viewer that is built and checks that each key equals the one
computed from this specification, by an independent implementation in the script, and that the
viewers' index files are byte-identical (\file{meta.json} compared without its informative
\code{engine} and \code{indexed\_s}):
\begin{lstlisting}[style=shell]
python tests/parity/index_parity.py                             # the synthetic test recordings
python tests/parity/index_parity.py --adapters-dir DIR FILE...  # recordings an adapter reads
\end{lstlisting}

\section{The files of a WFDB collection}
\label{sec:formats-reference-collection}

A WFDB collection is a folder of WFDB records with a few CSV files beside them that say how the
records fit together in time and what else the recording held: vital signs, alarms, breaks.
Converters write collections (the STP add-on's \code{wfdb} export is one), and anyone can write
one with a few lines of code. \file{index.csv} makes a folder a collection; the other files are
optional.
\begin{lstlisting}[style=text]
recording/                    the collection (one recording)
    index.csv                 the records, in order (required: it makes the folder a collection)
    gaps.csv                  breaks between records (optional)
    vitals.csv                vital-sign values (optional)
    events.csv                alarms and other events with a duration (optional)
    README.txt                anything for people (ignored by the readers)
    ecg/                      a folder per modality, each with its RECORDS file
        RECORDS
        ecg_250_0001.hea  ecg_250_0001.dat  ecg_250_0001.mrk
        ...
    pressure/
        ...
\end{lstlisting}

\begin{itemize}
\item \textbf{Records} are standard WFDB records (any signal format the WFDB library reads); each
  signal of each record becomes a piece of a channel. Every record's header states its sampling
  frequency: the readers assume none, so that a collection with a record that does not is not
  read. A signal with a gain of 0 (uncalibrated) is shown in stored counts, unit \code{adu}; a
  signal whose header gives no units has none.
\item \textbf{Records may sit in subfolders}; \file{index.csv} names them by their path relative
  to the collection, with \code{/} between folders (\code{ecg/ecg\_250\_0001}).
\item \textbf{CSV files} are UTF-8, comma-separated, with a header row, as Python's \code{csv}
  module writes them (quotes around the fields that need them). Columns not described here are
  ignored, so that a writer may add its own.
\item \textbf{Times} are Unix seconds (UTC), written as decimal numbers.
\end{itemize}

\subsection{index.csv}
\label{sec:formats-reference-indexcsv}

One row per record, in the order the records were recorded.

\begin{xltabular}{\linewidth}{@{}P{2.5cm}P{1.6cm}L@{}}
\caption{The columns of \file{index.csv}.}
\label{tab:formats-reference-indexcsv}\\
\toprule
\tabhead Column & \tabhead Required & \tabhead Meaning \\
\midrule
\endfirsthead
\toprule
\tabhead Column & \tabhead Required & \tabhead Meaning \\
\midrule
\endhead
\bottomrule
\endfoot
\code{record} & yes & the record's path relative to the collection, without extension; a name
  that is absolute or has a \code{..} component is not read (\cref{sec:recordings-collection}) \\
\code{modality} & & the kind of signals (\code{ecg}, \code{pressure}, \code{resp}, \code{spo2},
  \dots); channels are grouped by it. Default \code{wfdb} \\
\code{start\_unix} & & the time of the record's first sample. Empty: the record states no time
  (\cref{chap:recordings}) \\
\code{stretch} & & a label shared by the records of one stretch of contiguous data; records
  without a start time that share it are placed together \\
\end{xltabular}

Other columns (\code{fs}, \code{n\_sig}, \code{signals}, \code{units}, \code{sig\_len},
\code{duration\_s}, \code{start\_utc}, \code{time\_source}, \code{time\_uncertainty\_s},
provenance, checksums) are for people and other programs; the readers take the rate, the signals
and the length from the headers.

\subsection{gaps.csv}
\label{sec:formats-reference-gaps}

One row per break in the data, or per change that started a new record.

\begin{xltabular}{\linewidth}{@{}P{4.4cm}L@{}}
\caption{The columns of \file{gaps.csv}.}
\label{tab:formats-reference-gaps}\\
\toprule
\tabhead Column & \tabhead Meaning \\
\midrule
\endfirsthead
\toprule
\tabhead Column & \tabhead Meaning \\
\midrule
\endhead
\bottomrule
\endfoot
\code{scope} & \code{all} (every channel), or the modality or channels it concerns \\
\code{reason} & why: \code{signal\_loss}, \code{corrupt\_bytes}, \dots The reasons
  \code{channel\_absent}, \code{channel\_set\_change}, \code{label\_change} and
  \code{packet\_count\_change} mark where a record was split while the data stayed contiguous:
  they are not shown as breaks \\
\code{end\_before\_unix}, \code{start\_after\_unix} & the end of the data before and the start of
  the data after, when known \\
\code{gap\_s} & the length of the break, when known (even if neither end has a stated time) \\
\code{gap\_uncertainty\_s}, \code{measured\_by}, \code{exact}, \code{notes} & how it was measured:
  \code{exact} is \code{True} when \code{gap\_s} is exact \\
\code{records\_before}, \code{records\_after} & the records on each side, separated by \code{;},
  by name or path \\
\end{xltabular}

\subsection{vitals.csv}
\label{sec:formats-reference-vitals}

One row per value.

\begin{xltabular}{\linewidth}{@{}P{3.0cm}L@{}}
\caption{The columns of \file{vitals.csv}.}
\label{tab:formats-reference-vitalscsv}\\
\toprule
\tabhead Column & \tabhead Meaning \\
\midrule
\endfirsthead
\toprule
\tabhead Column & \tabhead Meaning \\
\midrule
\endhead
\bottomrule
\endfoot
\code{time} & when the value was measured \\
\code{label} (else \code{name}) & what was measured: \code{HR}, \code{SpO2}, \code{NBP-S}, \dots \\
\code{value} & the value (rows whose value is not a number are left out) \\
\code{unit} & its unit \\
\code{validity} & empty or \code{0}: valid; anything else: the row is left out \\
\code{source} & where the value came from, when a recording has several (the monitor's numerics,
  an import) \\
\end{xltabular}

\paragraph{Naming.} A label measured by one source is named by its label. When a label has values
from several sources, the first source in the file keeps the plain label and the others are named
with their source: \code{HR}, \code{HR [import]}. Series of different sources are never merged.

\subsection{events.csv}
\label{sec:formats-reference-eventscsv}

Alarms and other events that have a start and, possibly, an end.

\begin{xltabular}{\linewidth}{@{}P{2.6cm}L@{}}
\caption{The columns of \file{events.csv}.}
\label{tab:formats-reference-eventscsv}\\
\toprule
\tabhead Column & \tabhead Meaning \\
\midrule
\endfirsthead
\toprule
\tabhead Column & \tabhead Meaning \\
\midrule
\endhead
\bottomrule
\endfoot
\code{start\_unix} & the start (required) \\
\code{end\_unix} & the end; empty: an instant \\
\code{level} & the severity: 1 advisory, 2 warning, 3 crisis (shown by color; the readers' alarm
  colors are the only red on their pages) \\
\code{text} & what is shown \\
\code{kind} & the kind of event (\code{alarm}, the system that raised it, \dots); default
  \code{event} \\
\end{xltabular}

Events are kept when their start lies within a day of the recording's data (a guard against zero
dates), sorted by start.

\subsection{Markers}
\label{sec:formats-reference-markers}

A record may have a \file{.mrk} annotation file (WFDB's MIT format): marks at samples, such as R
waves or pacing spikes. The annotation's \code{chan} is the signal it belongs to. A marker's name
is its auxiliary text when it has one (\code{R\_WAVE}), else the WFDB mnemonic of its code
(\code{N}, \code{V}, \code{+}, \dots). Markers are drawn on their signal and can be listed by name.

\subsection{An example}
\label{sec:formats-reference-example}

The smallest collection: one ECG record with a start time, and an alarm.
\begin{lstlisting}[style=text]
index.csv
  record,modality,start_unix
  ecg/rec_0001,ecg,1709287200.0
events.csv
  start_unix,end_unix,level,text,kind
  1709287201,1709287203,2,HR HIGH 151 > 150,alarm
ecg/RECORDS
  rec_0001
ecg/rec_0001.hea, ecg/rec_0001.dat
\end{lstlisting}
\file{viewers/vitalstream/testdata/generate.py} writes a complete one (three stretches, a break of
stated length, a record without a start time, vital signs from two sources, alarms and markers)
with the Python standard library only; the three viewers and Medorama read it to the same index,
byte for byte.

\section{EDF, EDF+, BDF and BDF+: the reading rules}
\label{sec:formats-reference-edf}

The Rust engine (\file{libs/engine/src/sources/edf.rs}), the C++ engine
(\file{viewers/vitalstream/cpp/engine/src/sources\_edf.cpp}) and the Python viewer
(\file{viewers/vitalstream/python/vitalstream/vitalstream\_py/sources/edf.py}) implement the rules
of this section. Three implementations that read the same file must write byte-identical indexes
(\cref{sec:formats-reference-versions}). A 300\,MB EDF+ file of 28\,800 data records is indexed in
about 0.5 to 1 second by each of them. What users see is described in \cref{sec:recordings-edf}
and \cref{sec:vitalstream-edf}; the specifications of the formats are published at
\url{https://www.edfplus.info/specs/}.

\subsection{The header}
\label{sec:formats-reference-edf-header}

\begin{itemize}
\item The fixed part has 256 bytes. It is followed by 256 bytes per signal, in which each field
  is written for every signal in turn: labels, transducers, physical dimensions, physical minima,
  physical maxima, digital minima, digital maxima, prefilters, samples per data record, and the
  reserved field.
\item A field is its bytes read as Latin-1, with surrounding whitespace removed (as Python's
  \code{str.strip()} removes it).
\item The version field is \code{0} for EDF. For BDF, it is the byte \code{0xFF} followed by
  \code{BIOSEMI}. Anything else is refused (\code{not an EDF or BDF file}).
\item The fixed part's reserved field, once stripped, starts with \code{EDF+C}, \code{EDF+D},
  \code{BDF+C} or \code{BDF+D} for EDF+ and BDF+.
\item The number of header bytes must be $256 \times (\text{number of signals} + 1)$, with at
  least one signal.
\item Integer fields (the number of header bytes, the number of data records, the number of
  signals, and the samples per data record) consist of an optional sign and ASCII digits only.
\item Other numbers are ASCII decimal numbers: an optional sign, digits, a decimal point and an
  exponent, with surrounding whitespace and underscores between digits as Python's
  \code{float()} accepts them (signal-core's \code{parse\_float}). They must be finite.
\item A signal labelled \code{EDF Annotations} or \code{BDF Annotations} is an \emph{annotation
  signal}; every other signal is an \emph{ordinary signal}.
\item Signals with zero samples per data record are left out.
\item The duration of a data record must be positive when the file has an ordinary signal, and each
  ordinary signal's rate (its samples per data record divided by the duration) must be a finite
  number: a duration of \code{1e-320}\,s is refused (\code{signal \dots{} has no finite sampling
  rate}).
\item A file whose runs would make more than 100\,000 pieces (an EDF+D file of many short,
  separate data records) is refused with a message that says so
  (\cref{sec:formats-reference-limits}).
\item A \file{.edf} or \file{.bdf} file whose header cannot be read is still labelled \code{EDF}
  when browsing; it is refused when it is opened.
\end{itemize}

\subsection{The start time}
\label{sec:formats-reference-edf-start}

\begin{itemize}
\item The date (\code{dd.mm.yy}) and the time (\code{hh.mm.ss}) are 8-byte fields. Positions 0--1,
  3--4 and 6--7 are ASCII digits; positions 2 and 5 are any character that is not an ASCII digit.
  A non-ASCII separator makes the field invalid.
\item Two-digit years 85 to 99 are 1985 to 1999, and 00 to 84 are 2000 to 2084.
\item In EDF+ and BDF+, when the recording field starts with \code{Startdate } and its second
  word is 11 bytes long and ends in four ASCII digits, those four digits are the year. Only the
  year is used: the day and the month always come from the header's date, and the second word's
  dashes and month are not checked. The header's year positions are then ignored; any two
  characters are accepted there, typically \code{yy}, which EDF+ writes after 2084.
\item The month must be 1 to 12, the day 1 to 31 (not checked against the month: 31.02 becomes
  March 2), the hour 0 to 23, and the minute and second 0 to 59.
\item The start, in Unix seconds, is that date and time taken as UTC.
\item An invalid date or time means that the file states no start. Times then count from 0, and
  every piece is marked untimed (\code{timed: false}).
\end{itemize}

\subsection{Data records and physical values}
\label{sec:formats-reference-edf-records}

\begin{itemize}
\item The size of a data record is the sum of the samples per data record of all signals, times
  2 bytes (EDF) or 3 bytes (BDF).
\item Samples are little-endian two's-complement integers.
\item The number of records read is the number of whole records in the file after the header.
  When the header states a non-negative count smaller than that, only that many are read, and the
  records beyond it are left out without a note.
\item Bytes after the file's last whole record are left out, with a note.
\item A digital value $d$ becomes the physical value
  $p_{\min} + (d - d_{\min}) \times \bigl((p_{\max} - p_{\min}) / (d_{\max} - d_{\min})\bigr)$
  when $d_{\max} \neq d_{\min}$, and $d$ itself when $d_{\max} = d_{\min}$.
\item The value is computed in IEEE double precision, in exactly this order, without fused
  multiply-add. The envelopes store it as a 32-bit float.
\end{itemize}

\subsection{Annotations, onsets and runs}
\label{sec:formats-reference-edf-annotations}

\paragraph{Time-stamped annotation lists (TALs).}
\begin{itemize}
\item An annotation signal's bytes in a data record are TALs, each ended by a zero byte. Runs of
  zero bytes are padding.
\item A TAL is \code{onset[\textbackslash x15duration]\textbackslash x14text\textbackslash
  x14text\textbackslash x14\dots}.
\item The onset and the duration are numbers, read as above. A TAL whose onset is not valid UTF-8
  or not a number is skipped. A duration that is not a number counts as none, and the TAL is
  kept.
\item Each text is decoded as UTF-8, with invalid bytes replaced, and its runs of whitespace are
  collapsed to single spaces (as Python's \code{' '.join(text.split())} does).
\end{itemize}

\paragraph{Record onsets and runs.}
\begin{itemize}
\item In EDF+ and BDF+, a record's onset is the onset of the first TAL of the first annotation
  signal (its \emph{time-keeping annotation}), even when that signal has 0 samples per data
  record.
\item Otherwise, or when a record has no such TAL, a record starts where the record before it
  ends; the first one starts at 0.
\item A record whose onset differs by more than $10^{-6}$\,s from the end of the record before it
  starts a new \textbf{run}.
\item Each run makes one piece per ordinary signal, in signal order, with
  $t_0 = \text{start} + \text{onset of the run's first record}$.
\end{itemize}

\paragraph{Annotations.}
\begin{itemize}
\item Every non-empty text of every TAL is an event, with $t_0 = \text{start} + \text{onset}$,
  and $t_1 = t_0 + \text{duration}$ when the duration is positive, else $t_1 = t_0$. Its
  \code{level} is null and its \code{kind} is \code{annotation}.
\item The annotation signals of plain EDF and BDF files also give events; they only do not keep
  time.
\item Events are listed in the order of the file, then sorted by $t_0$ with a stable sort.
\end{itemize}

\subsection{The index of an EDF file}
\label{sec:formats-reference-edf-index}

The index follows \cref{sec:formats-reference-index}, with the following values.
\begin{itemize}
\item \code{kind} is \code{edf}; \code{format} is \code{EDF}, \code{EDF+}, \code{BDF} or
  \code{BDF+}.
\item \textbf{Channels.} There is one channel per ordinary signal, in signal order, created
  before any record is read. Each has a \code{label} (repeated labels get \code{ (2)},
  \code{ (3)} \dots), a \code{unit} (the physical dimension) and the \code{modality}
  \code{edf}.
\item \textbf{Pieces.} Besides the common keys (\code{id}, \code{ch}, \code{fs}, \code{n},
  \code{t0}, \code{timed}, \code{env}), each piece has \code{sig} (the signal's position in the
  header, annotation signals included) and \code{data\_record} (the run's first record), in that
  order.
\item \textbf{Gaps.} Between consecutive runs, when the next run starts after the previous one
  ends, there is a gap with \code{t0} (the end of the previous run), \code{t1}, \code{reason}
  \code{not recorded}, \code{gap\_s}, \code{exact} true, \code{measured\_by} null and
  \code{notes} empty.
\item \textbf{Alarms} are the events above.
\item \code{vitals} and \code{marker\_names} are empty, \code{skipped} is \code{\{\}}, and
  \code{reader} is \code{\{\}}.
\item \textbf{Notes}, in this order:
  \begin{enumerate}
  \item \code{The file's start time is taken as UTC.}, or \code{The file states no valid start
    time: times are counted from its start, and marked "time not stated".};
  \item then, when the header states more records than the file holds, \code{The header states N
    data records; the file holds M.}; otherwise, when bytes follow the last whole record,
    \code{The file ends inside a data record: its last N bytes are left out.}
  \end{enumerate}
\item A file with no ordinary signal, or with no whole data record, is refused.
\item \textbf{The key} is a file's key, from its size and modification time
  (\cref{sec:formats-reference-key}).
\end{itemize}

The header's patient and recording identification fields are never written into the index.
Only the year of an EDF+ recording field's \code{Startdate} is read, to date the recording.

\subsection{Test files}
\label{sec:formats-reference-edf-tests}

\file{viewers/vitalstream/testdata/edf/} holds synthetic files written by
\file{viewers/vitalstream/testdata/generate.py} (\cref{tab:formats-reference-edf-tests}). They are
checked by \file{libs/engine/tests/edf.rs}, by Medorama's \code{identify} and de-identification
tests, by the index parity test across the three viewers, and, for the ECG viewer, by
\file{viewers/ecg/tests/edf.test.js}.

\begin{xltabular}{\linewidth}{@{}P{3.0cm}L@{}}
\caption{The synthetic EDF test files.}
\label{tab:formats-reference-edf-tests}\\
\toprule
\tabhead File & \tabhead Tests \\
\midrule
\endfirsthead
\toprule
\tabhead File & \tabhead Tests \\
\midrule
\endhead
\bottomrule
\endfoot
\file{psg\_plus.edf} & EDF+C: three signals at 250, 25 and 1\,Hz; a first record 0.25\,s after
  the header's time; annotations with and without a duration; an invented identity in the
  patient field, which must appear in no index file. \\
\file{gappy\_plus.edf} & EDF+D: records of 0.5\,s in three runs, with breaks. \\
\file{plain.edf} & EDF of 1999 (a two-digit year); the number of records unknown (\code{-1}); a
  record cut short at the end. \\
\file{eeg\_plus.bdf} & BDF+C: 24-bit samples; a signal whose digital range is zero. \\
\end{xltabular}
